\section{Глава 6. Обучение на основе временных различий}

\subsection{Упражнение 6.1(c. 154)}
Если $V$ изменяется на протяжении эпизода, то $(6.6)$ справедливо только приближённо; какова разность между левой и правой частями? Обозначим $V_t$ массив ценностей состояний, используемый в выражении $(6.5)$ для TD-ошибки в момент $t$ и в правиле TD-обновления $(6.2)$. Повторив приведённый выше вывод, найдите дополнительный член, который нужно прибавить к сумме TD-ошибок, чтобы получить ошибку Монте-Карло.
\begin{equation}
    G_t - V(S_t) = \sum_{k=t}^{T-1} \gamma^{k-t} \delta_k.
    \tag{6.6}
\end{equation}
\begin{solution}
    Реально вычисляемая TD-ошибка использует версию массива на момент $k$:
    \[
    \delta_k = R_{k+1} + \gamma V_k(S_{k+1}) - V_k(S_k).
    \]
    Прибавим и вычтем $\gamma V_t(S_{t+1})$:
    \[
    G_t - V_t(S_t) = \delta_t + \gamma\left(G_{t+1} - V_t(S_{t+1})\right).
    \]
    Чтобы рекурсивно применить то же преобразование, нужно перейти от $V_t(S_{t+1})$ к $V_{t+1}(S_{t+1})$:
    \[
    G_{t+1} - V_t(S_{t+1}) = \left(G_{t+1} - V_{t+1}(S_{t+1})\right) + \left(V_{t+1}(S_{t+1}) - V_t(S_{t+1})\right).
    \]
    Раскручивая по индукции до конца эпизода (с учётом $G_T = V_T(S_T) = 0$), получаем:
    \[
    G_t - V_t(S_t) = \sum_{k=t}^{T-1} \gamma^{k-t} \delta_k + \sum_{k=t+1}^{T-1} \gamma^{k-t} \left(V_k(S_k) - V_{k-1}(S_k)\right).
    \]
    Таким образом, искомый дополнительный член равен
    \[
    \Delta_t = \sum_{k=t+1}^{T-1} \gamma^{k-t} \left(V_k(S_k) - V_{k-1}(S_k)\right). 
    \]
\end{solution}

\subsection{Упражнение 6.2(c.157)}
Это упражнение поможет развить интуитивное понимание того, почему TD-методы часто оказываются эффективнее методов Монте-Карло. Рассмотрим пример с поездкой домой и подход к его решению методами TD и Монте-Карло. Можете ли вы представить себе ситуацию, в которой TD-обновление в среднем было бы лучше обновления Монте-Карло? Приведите пример ситуации~-- описание прошлого опыта и текущего состояния,~-- в которой вы ожидаете, что TD-обновление будет лучше. 
\begin{solution}
    Пусть из старого офиса накоплен обширный опыт, и для всех состояний пути после въезда на трассу оценки $V(s)$ уже точны. После переезда меняется только начальный участок пути (новая парковка~$\to$ точка въезда на трассу), а хвост пути остаётся прежним.

    Для нового состояния <<новая парковка>> TD-обновление использует $V(S')$, где $S'$~-- точка въезда на трассу. Поскольку $V(S')$ уже точно отражает оставшееся время до дома, TD-таргет имеет малую дисперсию и сразу подтягивает к новому состоянию всю информацию, выученную из старого опыта. Достаточно нескольких поездок, чтобы оценка для новой парковки стала точной.

    Метод Монте-Карло использует фактически прошедшее время до дома~-- случайную величину, дисперсия которой накапливается по всему пути (пробки, светофоры, погода на знакомой части маршрута). Чтобы усреднить этот шум, требуется много эпизодов, хотя информация о хвосте пути уже была выучена ранее.

    Аналогичный эффект возникает и без переезда офиса: разные состояния пути встречаются с разной частотой, и редко посещаемые стартовые состояния обучаются медленно, тогда как промежуточные состояния хорошо изучены. TD переносит точную информацию из изученных состояний на соседние через одно обновление, что отражает использование марковской структуры задачи; метод Монте-Карло такого переноса не выполняет.
\end{solution}

\subsection{Упражнение 6.3(c. 159)}
Из результатов примера случайного блуждания, показанных на левом графике, может показаться, что первый эпизод приводит только к изменению $V(A)$. Какую информацию это несёт о том, что происходило в первом эпизоде? Почему изменилась только оценка для одного этого состояния? На сколько точно она изменилась?
\begin{solution}
    Эпизод всегда начинается в C, начальные ценности $V(s) = 0.5$ для всех неконечных состояний, шаг $\alpha = 0.1$, $\gamma = 1$. TD-ошибка одного шага:
    \[
    \delta_t = R_{t+1} + \gamma V(S_{t+1}) - V(S_t).
    \]
    При переходе между двумя неконечными состояниями $\delta_t = 0 + 0.5 - 0.5 = 0$, поэтому такие переходы не меняют оценок. Ненулевое обновление возможно лишь на последнем шаге — при переходе в терминал.

    Раз изменилась только $V(A)$, последний переход был из $A$ в левый терминал ($R = 0$, $V(\text{терм}) = 0$):
    \[
    \delta = 0 + 0 - 0.5 = -0.5.
    \]
    Тогда
    \[
    V(A) \leftarrow V(A) + \alpha \delta = 0.5 + 0.1 \cdot (-0.5) = 0.45.
    \]
    То есть эпизод завершился выходом влево из $A$, а $V(A)$ уменьшилась на $0.05$.
\end{solution}

\subsection{Упражнение 6.4(c. 159)}
Конкретные результаты, показанные на правом графике для примера случайного блуждания, зависят от размера шага $\alpha$. Как вы думаете, повлияло бы это на выводы о том, какой алгоритм лучше, если бы значения $\alpha$ изменялись в более широком диапазоне? Существует ли другое фиксированное значение $\alpha$, при котором любой из обсуждаемых алгоритмов работал бы значительно лучше, чем показано на рисунке? Объясните свой ответ.
\begin{solution}
    Качественные выводы не изменились бы: TD по-прежнему достигал бы меньшей асимптотической ошибки, чем MC. Это следствие свойств задачи, а не выбора $\alpha$.

    При любом постоянном $\alpha$ есть компромисс:
    \begin{itemize}
        \item большое $\alpha$ — быстрое убывание ошибки в начале, но высокий шум в пределе;
        \item малое $\alpha$ — медленное обучение, но низкая асимптотическая ошибка.
    \end{itemize}
    Расширение диапазона $\alpha$ в обе стороны добавило бы только худших кривых: при слишком больших $\alpha$ метод застревал бы на высоком уровне шума, при слишком малых — не успевал бы обучиться за 100 эпизодов.

    Значительно лучшего результата с фиксированным $\alpha$ добиться нельзя: для сходимости к истинным ценностям нужно убывающее $\alpha_n$, удовлетворяющее условиям:
    \[
    \sum_n \alpha_n = \infty, \quad \sum_n \alpha_n^2 < \infty.
    \]
    Только тогда метод и быстро движется в начале, и сходится в пределе.
\end{solution}

\subsection{*Упражнение 6.5(c. 159)}
На правом графике для примера случайного блуждания ошибка СКО TD-метода сначала убывает, а затем снова возрастает, особенно при больших значениях $\alpha$. Что могло бы быть причиной такого поведения? Как вы думаете, это происходит всегда или зависит от того, как инициализирована приближённая функция ценности?
\begin{solution}
    Причина — шум обновлений при большом постоянном $\alpha$. В стационарном режиме оценки не сходятся к истинным значениям, а «болтаются» вокруг них с некоторой амплитудой, определяемой самим $\alpha$ и динамикой задачи. Этот уровень шума не зависит от инициализации. Кривая сначала падает (обновления в среднем тянут к истинным значениям), достигает минимума, а затем подрастает до уровня шума.

    Сам эффект — стационарный шум выше минимума на траектории — присутствует при большом $\alpha$ \emph{всегда}, это свойство метода с постоянным шагом. Инициализация влияет лишь на то, насколько подъём визуально заметен: при старте близко к истинным значениям он виден, при старте далеко (например, $V(s) = 0$) маскируется большим начальным падением.
\end{solution}

\subsection{Упражнение 6.6(c. 159)}
В примере 6.2 мы сказали, что истинные ценности состояний от A до E в примере случайного блуждания равны $1/6, 2/6, 3/6, 4/6$ и $5/6$. Опишите по крайней мере два способа их вычисления. Как вы думаете, каким воспользовались мы? Почему?
\begin{solution}
    \textbf{Способ 1: уравнения Беллмана.} Политика равновероятна, $\gamma = 1$, ненулевая награда $+1$ возникает только при переходе из E в правый терминал. Уравнение Беллмана для $v_\pi$:
    \[
    v_\pi(s) = \sum_{a} \pi(a \mid s) \sum_{s', r} p(s', r \mid s, a)\bigl[r + \gamma v_\pi(s')\bigr].
    \]
    Подставляя политику и динамику задачи, получаем систему для пяти неконечных состояний:
    \begin{align*}
        v(A) &= \tfrac{1}{2} \cdot 0 + \tfrac{1}{2} v(B), \\
        v(B) &= \tfrac{1}{2} v(A) + \tfrac{1}{2} v(C), \\
        v(C) &= \tfrac{1}{2} v(B) + \tfrac{1}{2} v(D), \\
        v(D) &= \tfrac{1}{2} v(C) + \tfrac{1}{2} v(E), \\
        v(E) &= \tfrac{1}{2} v(D) + \tfrac{1}{2}(1 + 0).
    \end{align*}
    Каждое уравнение для внутреннего состояния говорит, что $v(s_i)$ равно среднему соседей, значит ценности линейны по номеру состояния: $v(s_i) = a + b \cdot i$. Граничные условия $v = 0$ при $i = 0$ и $v = 1$ при $i = 6$ дают $a = 0$, $b = 1/6$, откуда $v(s_i) = i/6$.

    \textbf{Способ 2: вероятностный аргумент.} Поскольку других вознаграждений нет и $\gamma = 1$, ценность состояния равна вероятности того, что эпизод, начатый из этого состояния, завершится в правом терминале. Это классическая задача о симметричном случайном блуждании с поглощающими барьерами, для которой известна формула $p(i) = i/N$ (здесь $N = 6$).

    \textbf{Каким способом воспользовались авторы.} Скорее всего вторым: знаменатель $6$ в ответе равен длине пути между терминалами и характерен именно для формулы случайного блуждания. Однако для того, кто не знаком с этим фактом, удобнее первый способ — он опирается только на уравнения Беллмана, введённые в этой книге, и приводит к тому же ответу через линейность решения и граничные условия.
\end{solution}

\subsection{*Упражнение 6.7(c. 162)}
Спроектируйте вариант обновления TD(0) с разделённой стратегией, который можно использовать с произвольной целевой стратегией $\pi$ и покрывающей поведенческой стратегией $b$ и в котором на каждом шаге $t$ применяется коэффициент выборки по значимости $\rho_{t:t}$ (5.3).
\begin{equation}
    \rho_{t:T-1} \doteq \frac{\prod_{k=t}^{T-1} \pi(A_k \mid S_k) p(S_{k+1} \mid S_k, A_k)}{\prod_{k=t}^{T-1} b(A_k \mid S_k) p(S_{k+1} \mid S_k, A_k)} = \prod_{k=t}^{T-1} \frac{\pi(A_k \mid S_k)}{b(A_k \mid S_k)}.
    \tag{5.3}
\end{equation}
\begin{solution}
    \leavevmode\\
    \begin{algorithm}[H]
        \DontPrintSemicolon
        \KwIn{целевая стратегия $\pi$, поведенческая стратегия $b$ (с покрытием $\pi$), шаг $\alpha \in (0, 1]$}
        Инициализировать $V(s)$ произвольно для всех $s \in \mathcal{S}^+$, $V(\text{terminal}) = 0$\;
        \ForEach{эпизод}{
            Инициализировать $S$\;
            \While{$S$ не терминальное}{
                $A \gets$ действие, выбранное согласно $b(\cdot \mid S)$\;
                Выполнить действие $A$, получить $R$ и $S'$\;
                $\rho \gets \dfrac{\pi(A \mid S)}{b(A \mid S)}$\;
                $V(S) \gets V(S) + \alpha \, \rho \, \bigl[R + \gamma V(S') - V(S)\bigr]$\;
                $S \gets S'$\;
            }
        }
        \caption{Off-policy TD(0) с одношаговым коэффициентом важности}
    \end{algorithm}
\end{solution}

\subsection{Упражнение 6.8(c. 163)}
Покажите, что вариант (6.6) для ценности действий имеет место, если ошибка TD для ценности действий определена как $\delta_t = R_{t+1} + \gamma Q(S_{t+1}, A_{t+1}) - Q(S_t, A_t)$ и снова предполагается, что ценности не изменяются от шага к шагу.
\begin{equation}
    G_t - V(S_t) = \sum_{k=t}^{T-1} \gamma^{k-t} \delta_k.
    \tag{6.6}
\end{equation}
\begin{solution}
    Замечание: в условии упражнения пропущен множитель $\gamma$ перед $Q(S_{t+1}, A_{t+1})$. Корректное определение TD-ошибки для ценности действий:
    \[
    \delta_t = R_{t+1} + \gamma Q(S_{t+1}, A_{t+1}) - Q(S_t, A_t).
    \]
    \begin{align*}
        G_t - Q(S_t, A_t)
        &= R_{t+1} + \gamma G_{t+1} - Q(S_t, A_t) \\
        &\quad + \gamma Q(S_{t+1}, A_{t+1}) - \gamma Q(S_{t+1}, A_{t+1}) \\
        &= \delta_t + \gamma\bigl(G_{t+1} - Q(S_{t+1}, A_{t+1})\bigr) \\
        &= \delta_t + \gamma \delta_{t+1}
        + \gamma^2 \bigl(G_{t+2} - Q(S_{t+2}, A_{t+2})\bigr) \\
        &= \delta_t + \gamma \delta_{t+1} + \gamma^2 \delta_{t+2} + \dots
        + \gamma^{T-t-1} \delta_{T-1} \\
        &\quad + \gamma^{T-t}\bigl(G_T - Q(S_T, A_T)\bigr) \\
        &= \delta_t + \gamma \delta_{t+1} + \gamma^2 \delta_{t+2} + \dots
        + \gamma^{T-t-1} \delta_{T-1} \\
        &\quad + \gamma^{T-t}(0 - 0) \\
        &= \sum_{k=t}^{T-1} \gamma^{k-t} \delta_k.
    \end{align*}
\end{solution}

\subsection{Упражнение 6.9. Ветреный сеточный мир с ходами короля(c. 164)}
Решите задачу о ветреном сеточном мире, считая, что возможных действий не четыре, а восемь, — добавляются ходы по диагонали. Насколько лучшее решение вы сможете найти, располагая дополнительными действиями? Можно ли ещё улучшить его, включив девятое действие — не двигаться самому, а положиться только на силу ветра?
\begin{solution}
    \textbf{Код смотри в отдельном файле репозитория.} \\
    Алгоритм Sarsa с параметрами $\varepsilon = 0.1$, $\alpha = 0.5$, $Q(s, a) = 0$ при инициализации. Число эпизодов увеличено до 2000 для надёжной сходимости в более сложных вариантах задачи.

    \textbf{Исходная задача (4 действия).} Минимальная длина жадной траектории — 15 шагов.
    \begin{figure}[H]
        \centering
        \includegraphics[width=0.7\textwidth]{original_trajectory.png}
        \caption{Жадная траектория для исходной задачи (15 шагов).}
    \end{figure}

    \textbf{Ходы короля (8 действий).} Диагональные ходы позволяют эффективнее компенсировать снос ветра. Минимальная длина жадной траектории — 7 шагов.
    \begin{figure}[H]
        \centering
        \includegraphics[width=0.7\textwidth]{diag_trajectory.png}
        \caption{Жадная траектория с ходами короля (7 шагов).}
    \end{figure}

    \textbf{Ходы короля + действие «стоять на месте» (9 действий).} Добавление девятого действия не улучшает результат: минимальная длина по-прежнему 7 шагов. Использовать ветер пассивно невыгодно, потому что бездействие тоже стоит $-1$ за шаг, а активное движение всегда даёт по крайней мере не худший прогресс.
    \begin{figure}[H]
        \centering
        \includegraphics[width=0.7\textwidth]{diag_stay_trajectory.png}
        \caption{Жадная траектория с ходами короля и действием «стоять на месте» (7 шагов).}
    \end{figure}

    \textbf{Вывод.} Переход от 4 к 8 действиям сократил оптимальную длину эпизода более чем вдвое (с 15 до 7 шагов). Девятое действие не приносит дополнительного выигрыша.
\end{solution}
\notebook{06_09_10}{6.9}

\subsection{Упражнение 6.10. Стохастический ветер(c. 164)}
Решите задачу о ветреном сеточном мире с ходами короля в предположении, что воздействие ветра если и есть, то стохастическое: иногда оно отклоняется на 1 от средних значений, указанных под столбцом. Точнее, треть времени сила ветра средняя (как в предыдущем упражнении), ещё треть времени вас сдувает на дополнительную ячейку вверх, а ещё треть — на дополнительную ячейку вниз.
\begin{solution}
    \textbf{Код смотри в отдельном файле репозитория.} \\
    Алгоритм Sarsa с теми же параметрами: $\varepsilon = 0.1$, $\alpha = 0.5$, $Q(s, a) = 0$, 2000 эпизодов. Действия — ходы короля (8 направлений). Сила ветра в столбце с базовым значением $w$ применяется как $w - 1$, $w$ или $w + 1$ с равными вероятностями $1/3$.

    \begin{figure}[H]
        \centering
        \includegraphics[width=0.7\textwidth]{6_10_trajectory.png}
        \caption{Одна из жадных траекторий при стохастическом ветре (12 шагов).}
    \end{figure}

    Из-за случайности ветра жадная траектория уже не детерминирована — каждый прогон даёт свою длину. По 100 прогонам жадной стратегии средняя длина эпизода составила $16.4 \pm 9.1$ шагов, с отдельными удачными прогонами в 12 шагов.

    \textbf{Вывод.} По сравнению с детерминированным случаем (7 шагов) стохастический ветер существенно усложняет задачу: оптимальная стратегия не может полагаться на точное знание сноса, и средняя длина эпизода более чем удваивается. Высокое стандартное отклонение отражает естественный разброс траекторий из-за случайности среды.
\end{solution}
\notebook{06_09_10}{6.10}

\subsection{Упражнение 6.11(c. 167)}
Почему Q-обучение считается методом управления с разделённой стратегией?
\begin{solution}
    В Q-обучении поведенческая и целевая политики различаются:
    \begin{itemize}
        \item \textbf{Поведенческая политика} (по ней выбираются действия в среде) — $\varepsilon$-жадная по $Q$, чтобы обеспечить покрытие;
        \item \textbf{Целевая политика} (та, чью функцию ценности мы оцениваем) — жадная по $Q$.
    \end{itemize}
    Это видно из самого правила обновления:
    \[
    Q(S, A) \leftarrow Q(S, A) + \alpha\Bigl[R + \gamma \max_{a'} Q(S', a') - Q(S, A)\Bigr].
    \]
    В качестве оценки следующего шага берётся $\max_{a'} Q(S', a')$ — ценность \emph{жадного} действия, а не реально выбранного $A'$ из $\varepsilon$-жадной политики. 
\end{solution}

\subsection{Упражнение 6.12(c. 167)}
Предположим, что выбор действия жадный. Будет ли тогда алгоритм Q-обучения в точности совпадать с Sarsa? Будут ли они выбирать в точности одинаковые действия и одинаково обновлять веса?
\begin{solution}
    \textbf{Обновления весов совпадают.} В Sarsa используется $Q(S', A')$, в Q-обучении — $\max_{a'} Q(S', a')$. Если $A'$ выбирается жадно, то $A' \in \arg\max_{a'} Q(S', a')$, и значит
    \[
    Q(S', A') = \max_{a'} Q(S', a').
    \]
    Поэтому при жадном выборе действий правила обновления тождественны по значению, независимо от того, как разрешаются ничьи.

    \textbf{Последовательности действий могут различаться.} В состояниях, где у нескольких действий одинаковое $Q$, жадный выбор разрешает ничью случайно. Если в двух алгоритмах ничьи разрешаются независимыми случайными механизмами, то выбранные действия могут отличаться. При фиксированном детерминированном правиле разрешения ничьих (одинаковом для обоих алгоритмов) последовательности действий совпадают.

    \textbf{Итог.} При жадном выборе действий обновления весов всегда идентичны. Последовательности действий идентичны при условии одинакового разрешения ничьих.
\end{solution}

\subsection{*Упражнение 6.13(c. 171)}
Как выглядят уравнения обновления для двойной версии алгоритма Expected Sarsa с $\varepsilon$-жадной целевой стратегией?
\begin{solution}
    Поддерживаем две независимые оценки $Q_1$ и $Q_2$. На каждом шаге с вероятностью $0.5$ обновляем одну из них, с вероятностью $0.5$ — другую. Политика, по которой берётся ожидание, строится из той же оценки, которую обновляем; ценности под суммой берутся из другой оценки. Это разделяет роли «выбор» и «оценка» и устраняет смещение максимизации.

    Обозначим $\varepsilon$-жадную политику относительно оценки $Q_i$ как
    \[
    \pi_i(a \mid s) =
    \begin{cases}
        1 - \varepsilon + \dfrac{\varepsilon}{|\mathcal{A}(s)|}, & \text{если } a = \arg\max_{a'} Q_i(s, a'), \\[6pt]
        \dfrac{\varepsilon}{|\mathcal{A}(s)|}, & \text{иначе}.
    \end{cases}
    \]

    Тогда правила обновления:
    \begin{itemize}
        \item С вероятностью $0.5$:
        \[
        Q_1(S, A) \leftarrow Q_1(S, A) + \alpha\Bigl[R + \gamma \sum_{a'} \pi_1(a' \mid S') \, Q_2(S', a') - Q_1(S, A)\Bigr].
        \]
        \item Иначе:
        \[
        Q_2(S, A) \leftarrow Q_2(S, A) + \alpha\Bigl[R + \gamma \sum_{a'} \pi_2(a' \mid S') \, Q_1(S', a') - Q_2(S, A)\Bigr].
        \]
    \end{itemize}
\end{solution}

\subsection{Упражнение 6.14(c. 173)}
Опишите, как можно переформулировать задачу Джека об аренде машин (пример 4.2) в терминах послесостояний. Почему в этой конкретной задаче такая постановка, скорее всего, приведёт к ускорению сходимости?
\begin{solution}
    \textbf{Послесостояние в задаче Джека.} Послесостояние — это конфигурация офисов сразу после ночного перегона машин, но до дневных запросов и возвратов. Если утром в офисах было $(s_1, s_2)$ и Джек перегнал $a$ машин, послесостояние равно $\sigma(s, a) = (s_1 - a, s_2 + a)$.

    \textbf{Почему это ускоряет сходимость.} Многие пары $(s, a)$ приводят к одному и тому же послесостоянию. Например, послесостояние $(8, 7)$ достигается из $(10, 5)$ при $a = -2$, из $(9, 6)$ при $a = -1$, из $(8, 7)$ при $a = 0$ и т.\,д. Дальнейшая динамика для всех этих пар идентична, поэтому различать их бессмысленно.

    В обычной формулировке через $q(s, a)$ алгоритм отдельно учит ценность каждой такой пары, разбивая статистику между ними. В формулировке через $v(\sigma)$ все эти пары обновляют одну общую оценку, что даёт:
    \begin{itemize}
        \item меньше обучаемых параметров (пространство послесостояний значительно меньше пространства пар $(s, a)$);
        \item больше данных на параметр — любой опыт из любой пары вносит вклад в одну и ту же оценку.
    \end{itemize}
    Оба эффекта ускоряют сходимость по сравнению со стандартной формулировкой.
\end{solution}
